% \documentclass[a4paper,12pt]{article}
\documentclass[a4paper,11pt]{article}
% -------------------------------------------------
% Set up the page margins.
% -------------------------------------------------
% \usepackage[text={6.5in,9.5in},centering]{geometry}
% \setlength{\topmargin}{-0.25in}
\usepackage[margin=1in]{geometry}

% -------------------------------------------------
% Load Packages here
% -------------------------------------------------
%\usepackage{hyperref}
\usepackage{graphicx,url,subfig}
\RequirePackage[OT1]{fontenc}
\RequirePackage{amsthm,amsmath,amssymb,amscd}
\RequirePackage[numbers]{natbib}
\RequirePackage[colorlinks,citecolor=blue,urlcolor=blue]{hyperref}
\usepackage{graphics}
\usepackage{enumerate}
\usepackage{datetime}
\usepackage{etex}
% \usepackage[notcite,notref,color]{showkeys}
\usepackage[normalem]{ulem} % either use this (simple) or
\usepackage{soul} % use this (many fancier options)
\makeatother
\numberwithin{equation}{section}
\allowdisplaybreaks
\usepackage{mathrsfs}


% -------------------------------------------------
% check references
% -------------------------------------------------
% \usepackage{refcheck}
% \usepackage[notref,notcite]{showkeys}
% \usepackage[notcite]{showkeys}
% \usepackage{showkeys}

% -------------------------------------------------
% Environments here
% -------------------------------------------------
\theoremstyle{plain}
\newtheorem{theorem}{Theorem}[section]
\newtheorem{corollary}[theorem]{Corollary}
\newtheorem{lemma}[theorem]{Lemma}
\newtheorem{proposition}[theorem]{Proposition}
\newtheorem{exercise}{Exercise}
\theoremstyle{definition}
\newtheorem{definition}[theorem]{Definition}
\newtheorem{hypothesis}[theorem]{Hypothesis}
\newtheorem{assumption}[theorem]{Assumption}

\newtheorem{remark}[theorem]{Remark}
\newtheorem{conjecture}[theorem]{Conjecture}
\newtheorem{example}[theorem]{Example}

% -------------------------------------------------
%  Define short hand symbols.
% -------------------------------------------------
\newcommand{\E}{\mathbb{E}}
\newcommand{\W}{\dot{W}}
\newcommand{\B}{\textbf{B}}
\newcommand{\D}{\mathbb{D}}
\newcommand{\ud}{\ensuremath{\mathrm{d}}}
\newcommand{\Ceil}[1]{\left\lceil #1 \right\rceil}
\newcommand{\Floor}[1]{\left\lfloor #1 \right\rfloor}
\newcommand{\sgn}{\text{sgn}}
\newcommand{\Lad}{\text{L}_{\text{ad}}^2}
\newcommand{\SI}[1]{\mathcal{I}\left[#1 \right]}
\newcommand{\SIB}[2]{\mathcal{I}_{#2}\left[#1 \right]}
\newcommand{\Indt}[1]{1_{\left\{#1 \right\}}}
\newcommand{\LadInPrd}[1]{\left\langle #1 \right\rangle_{\text{L}_\text{ad}^2}}
\newcommand{\LadNorm}[1]{\left|\left|  #1 \right|\right|_{\text{L}_\text{ad}^2}}
\newcommand{\Norm}[1]{\left\|  #1   \right\|}
\newcommand{\Ito}{It\^{o} }
\newcommand{\Itos}{It\^{o}'s }
\newcommand{\spt}[1]{\text{supp}\left(#1\right)}
\newcommand{\InPrd}[1]{\left\langle #1 \right\rangle}
\newcommand{\mr}{\textbf{r}}
\newcommand{\Ei}{\text{Ei}}
\newcommand{\arctanh}{\operatorname{arctanh}}
\newcommand{\ind}[1]{\mathbb{I}_{\left\{ {#1} \right\} }}

\DeclareMathOperator{\esssup}{\ensuremath{ess\,sup}}

\newcommand{\steps}[1]{\vskip 0.3cm \textbf{#1}}
\newcommand{\calB}{\mathcal{B}}
\newcommand{\calD}{\mathcal{D}}
\newcommand{\calE}{\mathcal{E}}
\newcommand{\calF}{\mathcal{F}}
\newcommand{\calG}{\mathcal{G}}
\newcommand{\calK}{\mathcal{K}}
\newcommand{\calH}{\mathcal{H}}
\newcommand{\calI}{\mathcal{I}}
\newcommand{\calL}{\mathcal{L}}
\newcommand{\calM}{\mathcal{M}}
\newcommand{\calN}{\mathcal{N}}
\newcommand{\calO}{\mathcal{O}}
\newcommand{\calT}{\mathcal{T}}
\newcommand{\calP}{\mathcal{P}}
\newcommand{\calR}{\mathcal{R}}
\newcommand{\calS}{\mathcal{S}}
\newcommand{\bbC}{\mathbb{C}}
\newcommand{\bbN}{\mathbb{N}}
\newcommand{\bbP}{\mathbb{P}}
\newcommand{\bbZ}{\mathbb{Z}}
\newcommand{\myVec}[1]{\overrightarrow{#1}}
\newcommand{\sincos}{\begin{array}{c} \cos \\ \sin \end{array}\!\!}
\newcommand{\CvBc}[1]{\left\{\:#1\:\right\}}
\newcommand*{\one}{{{\rm 1\mkern-1.5mu}\!{\rm I}}}
\newcommand{\RR}{\mathbb{R}}
\newcommand{\R}{\mathbb{R}}
\newcommand{\myEnd}{\hfill$\square$}
\newcommand{\ds}{\displaystyle}
\newcommand{\Shi}{\text{Shi}}
\newcommand{\Chi}{\text{Chi}}
\newcommand{\Daw}{\ensuremath{\mathrm{daw}}}
\newcommand{\Erf}{\ensuremath{\mathrm{erf}}}
\newcommand{\Erfi}{\ensuremath{\mathrm{erfi}}}
\newcommand{\Erfc}{\ensuremath{\mathrm{erfc}}}
\newcommand{\He}{\ensuremath{\mathrm{He}}}

\def\crtL{\theta}

\newcommand{\conRef}[1]{(#1)}

\DeclareMathOperator{\Lip}{\mathit{L}}
\DeclareMathOperator{\LIP}{Lip}
\DeclareMathOperator{\lip}{\mathit{l}}

\DeclareMathOperator{\Vip}{\overline{\varsigma}}
\DeclareMathOperator{\vip}{\underline{\varsigma}}
\DeclareMathOperator{\vv}{\varsigma}

% % Set up mdframe for questions.
\newcounter{NQ}
\newcounter{LQ}
\newcounter{ZQ}
\usepackage{xcolor}
\usepackage[framemethod=tikz]{mdframed}
\mdfsetup{%
	middlelinecolor=lightgray,
	middlelinewidth=2pt,
	backgroundcolor=red!20,
	roundcorner=10pt}
\newenvironment{NickQuestion}[1][] %
{\refstepcounter{NQ}
	\begin{mdframed}[backgroundcolor=lightgray]\textbf{Nick's Question \theNQ \: #1~~}~~~\noindent}%
	{\end{mdframed}}
\newenvironment{LeQuestion}[1][] %
{\refstepcounter{LQ}
	\begin{mdframed}[backgroundcolor=red!10]\textbf{Le's Question \theLQ \: #1~~}~~~\noindent}%
	{\end{mdframed}}
\newenvironment{ZhijianQuestion}[1][] %
{\refstepcounter{ZQ}
	\begin{mdframed}[backgroundcolor=blue!10]\textbf{Zhijian's Question \theZQ \: #1~~}~~~\noindent}%
	{\end{mdframed}}
\numberwithin{NQ}{section}
\numberwithin{LQ}{section}
\numberwithin{ZQ}{section}

\usepackage[notref,notcite]{showkeys}

\title{Nick Eisenberg's Meeting Notes}
\author{Le Chen, Nick Eisenberg, Zhijian Wu}
\date{October 2019}

\usepackage{natbib}
\usepackage{graphicx}


\begin{document}
	\begin{NickQuestion}[Details from Remark 1.1 page 3]
		Is this correct in showing that part (i) implies the existence of $K$. 
		
		Let $\psi \in C_0^\infty(\R^d)$ then we have
		\begin{align*}
			\int_{\R^d} \ud x \gamma(x) \psi(x) & = \int_{\R^d} \ud x \varphi(x) \mathcal{F}(\psi)(x)
			\\&= \int_{\R^d} \ud x \mathcal{F}(\varphi)(x) \psi(x)
			\\& = \int_{\R^d} \ud x \mathcal{F}(\sqrt{\varphi}\sqrt{\varphi})(x) \psi(x)
			\\& = \int_{\R^d} \ud x \bigg(\mathcal{F}(\sqrt{\varphi})*\mathcal{F}(\sqrt{\varphi})\bigg)(x) \psi(x)
		\end{align*}
		and since this is true for all test functions $\psi$ then we must have $\gamma (x) = \bigg(\mathcal{F}(\sqrt{\varphi})*\mathcal{F}(\sqrt{\varphi})\bigg)(x)$
	\end{NickQuestion}

\begin{NickQuestion}[top of page 6]
	Why does assumption A imply that $\alpha \in (0, 2\wedge d)$?
\end{NickQuestion}

\begin{NickQuestion}[page 8]
	Why do we have the representation (2.3)?
\end{NickQuestion}

\begin{NickQuestion}[Lemma 2.3]
	\begin{enumerate}
		\item Why does $f(cA) = c^{d-\alpha}f(A)$ imply that when $f$ is a density then $f(cx)=c^{-\alpha}f(x)$? 
	
		\item 	Does 
	\[
		\int_{\R^d}|\xi|^2|\mathcal{F}g(\xi)|^2 \ud \xi = \int_{\R^d} |\nabla g(x)|^2 \ud x ? 
	\]
	
		\item Why does $(\tilde{g}^2*f)(x)=B^{-\alpha}(g^2*f)(B^{-1}x)$. Here is what I have
			\begin{align*}
				B^{-\alpha}(g^2*f)(B^{-1}x) &= B^{-\alpha}\int_{\R^d}g^2(y)f\left(\frac{x}{B}-y\right) \ud y
				\\&= \int_{\R^d}g^2(y)f\left(x-By\right) \ud y
			\end{align*}
	since $B^{-\alpha}f\left(\frac{x}{B}-y\right) = f(x-By)$. Then, 
	\begin{align*}
		\int_{\R^d}g^2(y)f\left(x-By\right) \ud y &= B^{-d}\int_{\R^d}g^2\left(\frac{u}{B}\right)f(x-u) \ud u
		\\&= B^{-d/2}\int_{\R^d}\tilde{g}^2\left(u\right)f(x-u) \ud u
		\\& = B^{-d/2}(\tilde{g}^2*f)(x)
	\end{align*}
	and this shows that 
	\[
	B^{-d/2}(\tilde{g}^2*f)(x) = B^{-\alpha}(g^2*f)(B^{-1}x)
	\]
	
	\item If I am correct on (3) then 
	\[
	\langle \tilde{g}^2*f, \tilde{g}^2 \rangle = B^{d-\alpha} 	\langle g^2*f, g^2 \rangle 
	\]
	
	\item If I am incorrect in (3) then how do we use $(\tilde{g}^2*f)(x)=B^{-\alpha}(g^2*f)(B^{-1}x)$ to get that $\langle \tilde{g}^2*f, \tilde{g}^2 \rangle = B^{-\alpha} 	\langle g^2*f, g^2 \rangle $. We have that 
	\begin{align*}
		\langle \tilde{g}^2*f, \tilde{g}^2 \rangle &= \int_{\R^d}  (\tilde{g}^2*f)(x)\tilde{g}^2(x) \ud x
		\\&=  \int_{\R^d}  B^{-\alpha}(g^2*f)(B^{-1}x)\tilde{g}^2(x) \ud x
		\\&=  \int_{\R^d}  B^{-\alpha}(g^2*f)(B^{-1}x)B^{-d/2}g^2(B^{-1}x) \ud x
		\\&= B^{-\alpha - \frac{d}{2}} \int_{\R^d} (g^2*f)(B^{-1}x)g^2(B^{-1}x) \ud x
		\\&  = B^{-\alpha + \frac{d}{2}} \int_{\R^d} (g^2*f)(u)g^2(u) \ud u
		\\& = B^{-\alpha + \frac{d}{2}}\langle g^2*f, g^2 \rangle 
	\end{align*}
	and this is also different from the paper. 
\end{enumerate}
\end{NickQuestion}
	
\begin{NickQuestion}[page 11]
	When $d=3$ we know that $G(t,\cdot)$ is a measure so does $(G(t,\cdot)*p_\epsilon)(x) = \int p_\epsilon(x-y)G(t,\ud y)$?
\end{NickQuestion}

\begin{NickQuestion}[page 11 typos]
\begin{enumerate}

	\item on the middle of page you say $d \le 3$ right before (3.6) but I think it need to be $d=3$. 

	\item  When defining $L$ and $L_\epsilon$. There is a $\ud x$ in both integrals but there needs to be a $\ud t$. 
	
	\item When applying the change of variables to $L_n$ and $L_{n,\epsilon}$ the index on the sum I think should be $\Sigma_n$ not $S_n$. 
\end{enumerate}
\end{NickQuestion}

\begin{NickQuestion}
	How are the equations on the bottom of page 11 proved. For example, using (3.3), we see that 
	\[
		L_1(y_1)=\int_0^\infty e^{-t}\int_0^t G(s,y_1) \ud s \ud t
	\]
	but according to the equation for $L_n$ on the bottom of page 11 we should also have 
	\[
		L_1(y_1) = L(y_1)
	\]
	where 
	\[
		L(x) = \int_0^\infty e^{-t}G(t,x) \ud t
	\]
But it doesnt seem that $	L_1(y_1) = L(y_1)$. For example, consider $d=1$ then we have that $G(t,x) = \frac{1}{2}1_{|x| < t}$ and so 
\begin{align*}
		L_1(y_1)&=\int_0^\infty e^{-t}\int_0^t G(s,y_1) \ud s \ud t
		\\&= \frac{1}{2} \int_0^\infty e^{-t}\int_0^t 1_{\{|y_1| < s \}}\ud s \ud t
		\\& = \frac{1}{2}\int_0^\infty e^{-t}\int_{|y_1| \wedge t }^t \ud s \ud t
		\\& = \frac{1}{2}\int_0^\infty e^{-t}(t- (|y_1| \wedge t )) \ud t
		\\& = \frac{1}{2} \int_{|y_1|}^\infty e^{-t}(t- |y_1|  ) \ud t
\end{align*}
However 
\begin{align*}
L(y_1) &= \int_0^\infty e^{-t} G(t,y_1) \ud t
\\&= \frac{1}{2} \int_0^\infty e^{-t} 1_{\{|y_1|<t\}} \ud t
\\&= \frac{1}{2} \int_{|y_1|}^\infty e^{-t} \ud t
\end{align*}
and these do not seem to be equal. 
\end{NickQuestion}

\begin{NickQuestion}[Lemma 3.3]
	What are the details showing 
	\[
		\int_{(\R^d)^n} \left[\int_{(\R^d)^n}  \prod_{k=1}^n K(x_k-y_k)L_n(y_1,...,y_n) \ud y  \right]^2 \ud x = \int_{(\R^d)^n}|\mathcal{F}L_n(\xi_1,...,\xi_n)|^2 \mu(\ud \xi_1)...\mu(\ud \xi_n)
	\]
	
	for example suppose that $n=1$ so we want to show 
	\[
			\int_{\R^d} \left[\int_{\R^d} K(x-y)L_n(y) \ud y \right]^2 \ud x = \int_{\R^d}|\mathcal{F}L_n(\xi)|^2 \mu(\ud \xi)
	\]
	By Plancherels Theorem we have 
	\begin{align*}
			\int_{\R^d} \left[\int_{\R^d} K(x-y)L_n(y) \ud y \right]^2 \ud x = 	\int_{\R^d} \left[\mathcal{F}\left(\int_{\R^d} K(\cdot-y)L_n(y) \ud y\right)(\xi) \right]^2 \ud \xi
	\end{align*}
	However anyway I continued seemed not to work. I thought we should use the property that $\mathcal{F}(f)\mathcal{F}(g) = \mathcal{F}(f*g)$ and this would give us 
	\begin{align*}
	\int_{\R^d} \left[\mathcal{F}\left(\int_{\R^d} K(\cdot-y)L_n(y) \ud y\right)(\xi) \right]^2 \ud \xi & = \int_{\R^d} \mathcal{F}\left(\int_{\R^d} K(\cdot-y)L_n(y) \ud y * \int_{\R^d} K(\cdot-y)L_n(y) \ud y\right)(\xi)  \ud \xi
	\end{align*}
	but continuing from this didnt seem to work. 
\end{NickQuestion}

\begin{NickQuestion}[Lemma 3.4]
	\begin{enumerate}
		\item Why does $\gamma(y_k - y_k') = \int_{\R^d}K(y_k-x_k)K(y_k'-x_k)\ud x_k$. We know that $\gamma = K*K$ so 
		\begin{align*}
			\gamma(y_k - y_k') &= \int_{R^d} K(y_k-y_k'-x_k)K(x_k) \ud x_k
			\\&= \int_{R^d} K(y_k-u)K(u-x_k) \ud u
		\end{align*}
		by letting $u= y_k' +x_k $ 
		
		\item Where are the first and second equalities coming from?
	\end{enumerate}
\end{NickQuestion}

\begin{NickQuestion}[Lemma 3.6 page 14]
	How do we get this inequality right above equation (3.13)?
\end{NickQuestion}

\end{document}